%!TEX TS-program = xelatex
%!TEX encoding = UTF-8 Unicode
% Résumé LaTeX Template
%
% This template has been downloaded from:
% https://github.com/billryan/resume
%
% Author:
% Bin "Bill" Yuan <yuanbin2014@gmail.com>
%
% Adapted to be a Quarto template by Mitchell O'Hara-Wild
%
% Template license:
% CC BY-SA 4.0 (https://creativecommons.org/licenses/by-sa/4.0/)
%
% Everything before \begin{document}: class options, the accent colour, the
% personal information that fills the header, and Pandoc's own preamble.
%-------------------------------------------------------------------------------
% CONFIGURATIONS
%-------------------------------------------------------------------------------
$passoptions.latex()$
%
% A4 paper size by default, use 'letterpaper' for US letter
$doc-class.tex()$

% Configure page margins with geometry. The class has already loaded the
% package, so these are set with \geometry rather than as package options.
$if(geometry)$
\geometry{$for(geometry)$$geometry$$sep$,$endfor$}
$endif$

% LOCAL: no upstream equivalent -- upstream has no accent colour beyond black.
% `headcolor` is the same option name every other quarto-vitae template uses
% for its accent, applied here to the title rule and to every entry's
% flush-right date (see resume.cls). Left alone, the class's own black
% default stands.
$if(headcolor)$
\definecolor{accent}{HTML}{$headcolor$}
$endif$

%-------------------------------------------------------------------------------
%	PANDOC SUPPORT
%-------------------------------------------------------------------------------
% common.latex is required in every custom Quarto LaTeX template, and calls
% pandoc.tex, which in turn brings in syntax highlighting, tables.tex,
% graphics.tex, citations.tex, babel-lang.tex, tightlist.tex, biblio-config.tex
% and header-includes. Do not hand-copy any of those: they are Pandoc's to
% maintain, and a stale local copy silently drifts out of step with the reader.
$common.latex()$
$after-header-includes.latex()$
$hypersetup.latex()$

% Typography. font-settings.latex honours mainfont/sansfont/monofont and their
% *options and *fallback variants, so a user changes the résumé's fonts the
% same way they would in any other Quarto pdf. It runs after the class's own
% \setmainfont{TeX Gyre Termes} above, so a `mainfont` here wins. fonts.latex
% is NOT called: it would load lmodern and unicode-math over the top of what
% the class has already set up. \basicInfo sets the contact line in \sffamily
% regardless of mainfont, exactly as upstream does; give `sansfont` to change
% that line's face without touching the body.
$font-settings.latex()$

% Four further partials are deliberately not called here:
%   fonts.latex  see above.
%   before-title.tex, title.tex, before-body.tex, toc.tex  a résumé has no
%     title block or contents page; \name/\basicInfo below play that role, and
%     \section here writes no \addcontentsline, so a toc would be empty.
